\documentclass[11pt]{article}

\usepackage{amsmath, amssymb}
\usepackage{graphicx}
\usepackage{geometry}
\usepackage{hyperref}
\geometry{margin=1in}

\begin{document}

\begin{titlepage}
    \centering

    {\Large \textbf{A Unified Framework Field Theory for Drift, Coherence, and Regime Dynamics in Large Language Models}\par}
    \vspace{1.5cm}

    {\large \textbf{Nawder Loswin}\par}
    {\large TriadicFrameworks.org\par}
    \vspace{1cm}

    {\large 2026\par}
    \vfill

    \begin{abstract}
    Large Language Models (LLMs) exhibit systematic behaviors---drift branching, coherence waves, framework collisions, and regime-dependent reasoning---that remain poorly characterized in existing literature. Framework Field Theory (FFT) provides a structural model for these phenomena by treating reasoning as motion within a triadic substrate composed of operators, regimes, and coherence fields. This paper introduces FFT's core definitions, presents eight falsifiable predictions, demonstrates reproducible simulations using existing LLMs, outlines minimal engineering artifacts derived from the theory, and situates FFT within established traditions in information theory, systems theory, cognitive science, and AI alignment.
    \end{abstract}

    \vfill
\end{titlepage}

\end{document}
